% =====================================================================
%  LAPORAN HASIL PROYEK HPC
%  Akselerasi Exact k-Nearest Neighbors (k-NN) pada Asisten Tanya-Jawab
%  Dokumen Kepatuhan Internal (SIKA) Berbasis Retrieval-Augmented Generation
%
%  Kompilasi:  pdflatex laporan_hpc_knn.tex  (2x agar daftar isi terisi)
%  Butuh file: ugm-logo.png  (satu folder dengan .tex ini)
% =====================================================================

\documentclass[12pt,a4paper,oneside]{report}

% ---------- Paket dasar ----------
\usepackage[T1]{fontenc}
\usepackage[utf8]{inputenc}
\usepackage[english,indonesian]{babel}
\usepackage{mathptmx}              % Times New Roman
\usepackage[a4paper,left=4cm,right=3cm,top=3cm,bottom=3cm]{geometry}
\usepackage{graphicx}
\usepackage{amsmath,amssymb}
\usepackage{booktabs}
\usepackage{longtable}
\usepackage{tabularx}
\usepackage{array}
\usepackage{enumitem}
\usepackage{setspace}
\usepackage{xcolor}
\usepackage{caption}
\usepackage{listings}
\usepackage{fancyhdr}
\usepackage{titlesec}
\usepackage{float}
\usepackage{pgfplots}
\usepackage[hidelinks]{hyperref}
\usepackage{tocloft}
\pgfplotsset{compat=1.17}

% ---------- Daftar isi/tabel/gambar kompak ----------
\renewcommand{\cftchapfont}{\normalfont\bfseries}
\renewcommand{\cftchappagefont}{\normalfont\bfseries}
\renewcommand{\cftsecfont}{\normalfont}
\renewcommand{\cftsubsecfont}{\normalfont}
\setlength{\cftbeforechapskip}{0.55em}
\setlength{\cftbeforesecskip}{0.05em}
\setlength{\cftbeforesubsecskip}{0em}
\setlength{\cftbeforetabskip}{0.3em}
\setlength{\cftbeforefigskip}{0.3em}
\setlength{\cftchapindent}{0em}
\setlength{\cftchapnumwidth}{3em}
\setlength{\cftsecindent}{2em}
\setlength{\cftsubsecindent}{4em}
\setlength{\cfttabindent}{0em}
\setlength{\cftfigindent}{0em}
\renewcommand{\cfttoctitlefont}{\hfill\large\bfseries}
\renewcommand{\cftaftertoctitle}{\hfill}
\renewcommand{\cftlottitlefont}{\hfill\large\bfseries}
\renewcommand{\cftafterlottitle}{\hfill}
\renewcommand{\cftloftitlefont}{\hfill\large\bfseries}
\renewcommand{\cftafterloftitle}{\hfill}

% ---------- Pengaturan umum ----------
\onehalfspacing
\setlength{\parindent}{1.5em}
\setlength{\parskip}{0.25em}
\renewcommand{\arraystretch}{1.3}

% ---------- Format bab & subbab ----------
\addto\captionsindonesian{%
  \renewcommand{\chaptername}{BAB}%
  \renewcommand{\contentsname}{DAFTAR ISI}%
  \renewcommand{\listtablename}{DAFTAR TABEL}%
  \renewcommand{\listfigurename}{DAFTAR GAMBAR}%
}
\renewcommand{\chaptername}{BAB}
\renewcommand{\thechapter}{\Roman{chapter}}
\renewcommand{\thesection}{\arabic{chapter}.\arabic{section}}
\renewcommand{\thesubsection}{\thesection.\arabic{subsection}}
\renewcommand{\thesubsubsection}{\thesubsection.\arabic{subsubsection}}
\renewcommand{\thetable}{\arabic{chapter}.\arabic{table}}
\renewcommand{\thefigure}{\arabic{chapter}.\arabic{figure}}
\renewcommand{\theequation}{\arabic{chapter}.\arabic{equation}}

\titleformat{\chapter}[display]
  {\normalfont\large\bfseries\centering}
  {\MakeUppercase{\chaptername}\ \thechapter}
  {0.5em}
  {\large\MakeUppercase}
\titlespacing*{\chapter}{0pt}{-1.5em}{1.5em}

\titleformat{\section}{\normalfont\large\bfseries}{\thesection}{0.8em}{}
\titleformat{\subsection}{\normalfont\normalsize\bfseries}{\thesubsection}{0.7em}{}
\titleformat{\subsubsection}{\normalfont\normalsize\bfseries}{\thesubsubsection}{0.6em}{}

% ---------- Header / footer ----------
\setlength{\headheight}{14pt}
\pagestyle{fancy}
\fancyhf{}
\fancyhead[L]{\small\itshape Laporan Proyek HPC}
\fancyhead[R]{\small\itshape Exact k-NN Vector Search (SIKA)}
\fancyfoot[C]{\thepage}
\renewcommand{\headrulewidth}{0.4pt}
\fancypagestyle{plain}{%
  \fancyhf{}%
  \fancyfoot[C]{\thepage}%
  \renewcommand{\headrulewidth}{0pt}%
}

% ---------- Caption ----------
\captionsetup{font=small,labelfont=bf,justification=centering}

% ---------- Listing gaya ----------
\definecolor{codebg}{gray}{0.96}
\lstdefinestyle{cppstyle}{
  backgroundcolor=\color{codebg},
  basicstyle=\ttfamily\footnotesize,
  keywordstyle=\bfseries,
  commentstyle=\itshape,
  numbers=left,
  numberstyle=\tiny,
  frame=single,
  rulecolor=\color{gray!50},
  breaklines=true,
  showstringspaces=false,
  tabsize=4,
  captionpos=b
}

% ---------- Perintah bantu ----------
\newcommand{\keyword}[1]{\textbf{\textit{#1}}}
\newcommand{\vect}[1]{\mathbf{#1}}
\newcommand{\pass}{\textcolor{green!45!black}{PASS}}

% toleransi pemenggalan kata (banyak istilah teknis panjang)
\setlength{\emergencystretch}{3em}
\tolerance=2000
\hbadness=10000

% =====================================================================
\begin{document}
\sloppy
\pagenumbering{roman}

% =====================================================================
%  HALAMAN SAMPUL
% =====================================================================
\begin{titlepage}
\centering
\singlespacing
\thispagestyle{empty}
\pagenumbering{alph}
\setlength{\parskip}{0pt}

\vspace*{0.2cm}

{\large\bfseries LAPORAN HASIL PROYEK\\[0.15em]
MATA KULIAH ARSITEKTUR DAN INFRASTRUKTUR BERKINERJA TINGGI\par}

\vspace{0.7cm}

\rule{\linewidth}{1pt}\par
\vspace{0.4em}
{\Large\bfseries AKSELERASI \MakeUppercase{Exact} $k$-NEAREST NEIGHBORS ($k$-NN)\\[0.2em]
PADA ASISTEN TANYA-JAWAB DOKUMEN KEPATUHAN INTERNAL\\[0.2em]
BERBASIS RETRIEVAL-AUGMENTED GENERATION (SIKA)\par}
\vspace{0.4em}
\rule{\linewidth}{1pt}\par

\vspace{0.6cm}
\includegraphics[width=6cm]{ugm-logo.png}\par

\vspace{0.6cm}
{\normalsize\bfseries Disusun oleh:\par}
\vspace{0.2em}

{\normalsize
\renewcommand{\arraystretch}{1.25}
\begin{tabular}{l l}
\textbf{Nama} & \textbf{NIM} \\[0.2em]
Axelle Chandra            & 24/533796/PA/22614 \\
Indratanaya Budiman       & 24/534784/PA/22683 \\
Kevin Antonio Wiyono Lauw & 24/535917/PA/22736 \\
Ayasha Rahmadinni         & 24/545462/PA/23178 \\
\end{tabular}\par}

\vfill

\vspace{0.3cm}

\end{titlepage}

% =====================================================================
%  DAFTAR ISI / TABEL / GAMBAR
% =====================================================================
\pagenumbering{roman}
\setcounter{page}{1}
\pagestyle{plain}

\begingroup
\setlength{\parskip}{0pt}
\renewcommand{\addvspace}[1]{}
\tableofcontents
\endgroup

\clearpage
\begingroup
\setlength{\parskip}{0pt}
\listoftables
\endgroup

\clearpage
\begingroup
\setlength{\parskip}{0pt}
\listoffigures
\endgroup

\clearpage
\pagestyle{fancy}
\pagenumbering{arabic}
\setcounter{page}{1}

% =====================================================================
\chapter{PENDAHULUAN}
% =====================================================================

\section{Latar Belakang Masalah}

Sistem \keyword{Retrieval-Augmented Generation} (RAG) menjawab pertanyaan dengan lebih dulu mengambil potongan dokumen yang paling relevan dari sebuah basis pengetahuan, lalu menyusun jawaban berdasarkan konteks tersebut. Pada domain kepatuhan internal, proses pengambilan ini tidak boleh melewatkan pasal pengecualian atau ambang batas nominal. Karena alasan itu, proyek ini memilih \keyword{Exact} $k$-Nearest Neighbors ($k$-NN) dan bukan \keyword{Approximate Nearest Neighbor} (ANN), agar \keyword{recall} tetap 100\%.

Perhitungan \keyword{exact} $k$-NN secara sekuensial pada CPU memiliki kompleksitas waktu $\mathcal{O}(Q \times N \times D)$, dengan $Q$ jumlah kueri, $N$ jumlah vektor basis data, dan $D$ jumlah dimensi. Pada skala ratusan ribu vektor, pendekatan sekuensial menjadi lambat dan tidak memenuhi standar latensi \keyword{real-time}. Ada dua \keyword{bottleneck} utama pada arsitektur komputer:

\begin{enumerate}[label=(\arabic*),leftmargin=2.2em]
  \item \keyword{Memory-bound}, yaitu matriks basis data berukuran ratusan megabita tidak muat di \keyword{L3 Cache} CPU sehingga terjadi \keyword{cache thrashing}.
  \item \keyword{Compute-bound}, yaitu instruksi aritmatika belum memanfaatkan paralelisasi data (SIMD) secara optimal.
\end{enumerate}

Proyek ini menjawab kedua masalah tersebut dengan pendekatan \keyword{High-Performance Computing} (HPC). Implementasi dan analisis komparatif dilakukan untuk algoritma \keyword{Exact} $k$-NN menggunakan paralelisasi \keyword{shared-memory} (OpenMP), akselerasi GPU (CUDA), dan distribusi memori antar-simpul (MPI). Optimasi diarahkan pada satu permasalahan operasional yang konkret, yaitu lamanya penelusuran dokumen kepatuhan internal pada aplikasi \keyword{SIKA} (Sistem Informasi Kepatuhan Internal).

\section{Rumusan Masalah}

Berdasarkan latar belakang di atas, rumusan masalah dalam proyek ini adalah:
\begin{enumerate}[leftmargin=2em]
  \item Bagaimana karakteristik \keyword{bottleneck} performa (latensi, \keyword{cache miss rate}, dan \keyword{instruction throughput}) dari algoritma \keyword{Exact} $k$-NN sekuensial pada data berdimensi tinggi?
  \item Bagaimana memetakan dan mengoptimasi operasi jarak vektor memanfaatkan hierarki memori (\keyword{loop tiling} di CPU dan \keyword{shared memory} di GPU)?
  \item Bagaimana merancang distribusi data (\keyword{sharding}) dan komunikasi antar-simpul dengan MPI untuk memproses data yang melebihi kapasitas memori satu proses?
  \item Seberapa besar peningkatan performa (\keyword{speedup}) dan efisiensi skalabilitas yang dihasilkan OpenMP, CUDA, dan MPI dibandingkan \keyword{baseline}?
  \item Sejauh mana optimasi tersebut membawa \keyword{retrieval} pada SIKA memenuhi target latensi p95 $\leq 100$~ms, dan bagaimana wujud sistem yang mendemonstrasikannya?
\end{enumerate}

\section{Tujuan Proyek}

Tujuan dari pelaksanaan proyek ini adalah:
\begin{enumerate}[leftmargin=2em]
  \item Mengembangkan dan memprofil program \keyword{baseline} sekuensial di CPU untuk mengukur latensi komputasi dan masalah \keyword{cache}.
  \item Mengimplementasikan paralelisasi \keyword{loop} dan vektorisasi (SIMD) dengan OpenMP.
  \item Mengembangkan \keyword{kernel} CUDA berbasis GEMM dengan \keyword{shared memory} dan \keyword{parallel reduction}.
  \item Menerapkan pola \keyword{sharding} dan \keyword{Gather} dengan MPI untuk menurunkan jejak memori per proses.
  \item Melakukan \keyword{benchmark} dan analisis metrik kinerja untuk membuktikan efektivitas setiap level arsitektur, dengan kriteria kelulusan yang ditetapkan sebelum pengukuran.
  \item Mengintegrasikan \keyword{engine} hasil optimasi ke dalam aplikasi SIKA yang dapat didemonstrasikan langsung, dengan \keyword{engine} yang dapat ditukar saat berjalan.
\end{enumerate}

\section{Batasan Masalah}

\begin{enumerate}[leftmargin=2em]
  \item Algoritma difokuskan pada \keyword{Brute-Force / Exact Search} $k$-NN dengan \keyword{Euclidean Distance}. Proyek tidak mencakup ANN (HNSW, FAISS) dan tidak membandingkan akurasi \keyword{Exact} terhadap ANN.
  \item Tipe data vektor adalah \keyword{single-precision floating-point} (\texttt{float32}).
  \item Dataset berupa matriks vektor berdimensi statis ($D=384$ atau $768$), sintetis maupun subset dataset publik, dengan jumlah vektor maksimum $500.000$.
  \item Skenario aplikatif diwujudkan sebagai purwarupa \keyword{single-node}. Aspek non-fungsional sistem produksi (autentikasi, ketersediaan tinggi, replikasi) di luar cakupan.
  \item Kualitas model LLM dan model \keyword{embedding} tidak dievaluasi dan diperlakukan sebagai komponen tetap.
  \item Evaluasi performa dilakukan pada lingkungan spesifik yang tersedia, didefinisikan pada Bab III.
\end{enumerate}

% =====================================================================
\chapter{TINJAUAN PUSTAKA DAN LANDASAN TEORI}
% =====================================================================

\section{Konsep \texorpdfstring{$k$}{k}-Nearest Neighbors pada Vector Database}

Sistem RAG mengandalkan pencarian informasi semantik melalui representasi vektor berdimensi tinggi (\keyword{embeddings}). Untuk menemukan dokumen yang relevan dengan kueri, sistem menjalankan algoritma $k$-NN untuk mencari $k$ vektor dalam basis data yang jaraknya paling dekat dengan vektor kueri. Untuk vektor kueri $\vect{q}$ dan vektor dokumen $\vect{p}$ yang berdimensi $D$, kuadrat jarak Euclidean dirumuskan sebagai:
\begin{equation}
d^{2}(\vect{q}, \vect{p}) = \sum_{i=1}^{D} (q_i - p_i)^{2}
\end{equation}

Pada komputasi berbasis matriks dan GPU, persamaan di atas diekspansi secara aljabar menjadi:
\begin{equation}
d^{2}(\vect{q}, \vect{p}) = \sum_{i=1}^{D} q_i^{2} + \sum_{i=1}^{D} p_i^{2} - 2 \sum_{i=1}^{D} q_i p_i
\end{equation}
atau dalam notasi aljabar linear:
\begin{equation}
d^{2}(\vect{q}, \vect{p}) = \lVert \vect{q} \rVert^{2} + \lVert \vect{p} \rVert^{2} - 2(\vect{q} \cdot \vect{p})
\end{equation}

Suku $2(\vect{q} \cdot \vect{p})$ dapat dipetakan menjadi operasi \keyword{General Matrix-Matrix Multiplication} (GEMM). Operasi ini sangat teroptimasi pada \keyword{hardware}, baik di CPU maupun GPU. Perlu dicatat bahwa ekspansi ini rentan \keyword{catastrophic cancellation} pada \keyword{float32} karena nilai $\lVert \vect{q} \rVert^{2} + \lVert \vect{p} \rVert^{2}$ jauh lebih besar daripada $d^{2}$. Bagian III menjelaskan cara proyek ini menjaga ketepatan hasil.

\section{Hierarki Memori dan \texorpdfstring{\textit{Cache Locality}}{Cache Locality}}

Performa algoritma tidak hanya bergantung pada jumlah instruksi aritmatika, tetapi juga pada kecepatan suplai data ke unit pemrosesan. Arsitektur komputer modern memakai hierarki memori:
\begin{enumerate}[leftmargin=2em]
  \item \keyword{L1 dan L2 Cache}, berukuran puluhan hingga ratusan Kilobita, bersifat \keyword{private} untuk setiap \keyword{core}, dengan latensi beberapa siklus \keyword{clock}.
  \item \keyword{L3 Cache (Last Level Cache)}, berukuran belasan hingga puluhan Megabita dan dipakai bersama seluruh \keyword{core} dalam satu cip.
  \item \keyword{Main Memory} (RAM), berkapasitas besar namun berlatensi ratusan siklus \keyword{clock}.
\end{enumerate}

Pada implementasi sekuensial yang naif, membaca matriks basis data berulang kali menyebabkan data di L3 Cache terusir sebelum sempat dipakai ulang. Fenomena ini disebut \keyword{Cache Thrashing}. Optimasi HPC memanfaatkan prinsip \keyword{Locality of Reference}: \keyword{temporal locality} melalui \keyword{Loop Tiling} (cache blocking), dan \keyword{spatial locality} melalui pembacaan data yang berdekatan secara sekuensial agar sejalan dengan \keyword{Cache Line}.

\section{Hukum Amdahl}

Dalam komputasi paralel, \keyword{speedup} dibatasi oleh porsi program yang berjalan sekuensial. Hukum Amdahl merumuskan batas ini:
\begin{equation}
S = \frac{1}{(1 - P) + \frac{P}{N}}
\end{equation}
dengan $S$ \keyword{speedup} teoretis, $P$ proporsi algoritma yang dapat diparalelkan, $1-P$ proporsi sekuensial, dan $N$ jumlah \keyword{thread}. Hukum ini menjadi dasar bahwa penambahan \keyword{core} tidak akan memberi peningkatan linier tanpa batas selama \keyword{overhead} sinkronisasi dan bagian sekuensial (misalnya pengurutan hasil akhir) masih dominan.

\section{Paradigma Pemrosesan Paralel}

\subsection{OpenMP (Shared-Memory Multiprocessing)}
OpenMP adalah API untuk \keyword{shared-memory multiprocessing} di C/C++. Beban \keyword{loop} dibagi ke beberapa \keyword{thread}. Arahan SIMD (\keyword{Single Instruction, Multiple Data}) memungkinkan kompilator memakai unit vektor perangkat keras (misalnya AVX2) untuk memproses beberapa elemen dimensi dalam satu siklus.

\subsection{CUDA (GPU Acceleration)}
CUDA adalah platform komputasi paralel dari NVIDIA. GPU diarsiteki dengan ribuan \keyword{core} kecil untuk eksekusi masif dengan \keyword{throughput} tinggi. Pada optimasi $k$-NN, pemanfaatan \keyword{library} aljabar linear (cuBLAS) dan \keyword{shared memory} menjadi kunci untuk meminimalkan ketergantungan \keyword{bandwidth} memori global.

\subsection{MPI (Distributed-Memory Message Passing)}
Ketika ukuran matriks melampaui kapasitas RAM satu \keyword{node}, paradigma terdistribusi diperlukan. MPI adalah standar untuk membagi komputasi lintas proses yang berkomunikasi melalui jaringan. Pola yang dipakai adalah \keyword{broadcast} untuk data kecil (vektor kueri) dan \keyword{Gather} untuk mengumpulkan \keyword{top-k} parsial ke proses utama untuk pengurutan final.

% =====================================================================
\chapter{METODOLOGI DAN SKENARIO IMPLEMENTASI}
% =====================================================================

\section{Lingkungan Pengujian}

Pengujian dilakukan pada lingkungan \keyword{single-node} agar dapat direproduksi dengan perangkat yang tersedia.

\begin{table}[H]
\centering
\caption{Spesifikasi lingkungan pengujian.}
\label{tab:env}
\begin{tabularx}{\linewidth}{@{}l X@{}}
\toprule
\textbf{Komponen} & \textbf{Spesifikasi}\\
\midrule
CPU & Intel Core i3-12100F, 4 core fisik / 8 thread, AVX2; L3 Cache 12~MB\\
Memori & RAM 16~GB DDR4 (WSL 8~GB)\\
GPU & NVIDIA GTX 1650, 4~GB VRAM, \keyword{compute capability} 7.5\\
Sistem Operasi & Windows 11 dengan WSL2 Ubuntu 24.04 LTS\\
Kompilator & GCC 13.3 (\texttt{-O3 -march=native -fopenmp})\\
CUDA & NVIDIA CUDA Toolkit 12.0 dengan cuBLAS\\
MPI & OpenMPI 4.x, 1--4 proses pada satu mesin (\keyword{oversubscribe})\\
\bottomrule
\end{tabularx}
\end{table}

Konfigurasi CPU dengan 4 \keyword{core} fisik dan LLC 12~MB membuat \keyword{bottleneck} memori bukan asumsi teoretis. Matriks basis terbesar ($500.000 \times 384 \times 4$ byte $\approx 768$~MB) melampaui LLC lebih dari 60 kali, sehingga selisih performa sebelum dan sesudah optimasi akan terlihat jelas pada perangkat yang tidak istimewa.

\section{Desain Algoritma Sekuensial (Baseline)}

Program \keyword{baseline} adalah CPU \keyword{single-thread} tanpa optimasi. Alur kerjanya: memuat matriks kueri dan basis ke memori, menjalankan dua \keyword{nested loop} (luar atas kueri, dalam atas basis), menghitung kuadrat jarak Euclidean tiap pasangan, lalu menyimpan \keyword{top-k} dengan \keyword{bounded max-heap} berukuran $k$. Pengukuran memakai pustaka \texttt{chrono} beresolusi tinggi. Selain mode \keyword{batch}, program menyediakan mode \keyword{single-query} yang mencatat waktu tiap kueri untuk menghasilkan nilai p50, p95, dan p99.

\section{Skenario Optimasi 1: Shared-Memory (OpenMP)}

Pendekatan OpenMP diterapkan dalam tiga tahap:
\begin{enumerate}[leftmargin=2em]
  \item \keyword{Loop-level parallelism}: \texttt{\#pragma omp parallel for schedule(dynamic)} pada blok kueri, agar beban terbagi dan ketidakseimbangan berkurang.
  \item \keyword{Vektorisasi SIMD}: \texttt{\#pragma omp simd} pada \keyword{loop} dimensi agar unit vektor (AVX) memproses beberapa elemen sekaligus.
  \item \keyword{Cache blocking} (\keyword{loop tiling}): baris basis dibaca sekali lalu dipakai ulang untuk sekumpulan kueri, sehingga jumlah pembacaan basis turun dari $Q$ kali menjadi $Q/\text{blok}$ kali.
\end{enumerate}

\section{Skenario Optimasi 2: GPU Acceleration (CUDA)}

Implementasi CUDA menggeser \keyword{nested loop} ke perangkat grafis:
\begin{enumerate}[leftmargin=2em]
  \item \keyword{Pemetaan memori}: matriks basis dipindahkan dari \keyword{host} ke VRAM sekali saat \keyword{engine} diaktifkan, bukan per kueri.
  \item \keyword{Konversi GEMM}: komputasi $\vect{q} \cdot \vect{p}$ dipetakan menjadi perkalian matriks $\mathbf{Q} \cdot \mathbf{B}^{\top}$ melalui cuBLAS.
  \item \keyword{Penyempurnaan eksak dan reduksi}: kandidat teratas hasil GEMM dihitung ulang dengan rumus eksak $\sum (q_i - p_i)^2$ untuk menjamin hasil identik \keyword{baseline}, lalu \keyword{top-k} dipilih lewat reduksi di \keyword{shared memory}.
\end{enumerate}

Langkah ketiga diperlukan karena ekspansi GEMM kehilangan presisi pada \keyword{float32}. Tanpa penyempurnaan, urutan tetangga dapat bergeser dari hasil \keyword{baseline}. Dengan penyempurnaan, kriteria kesamaan himpunan hasil (K5) terpenuhi tanpa mengorbankan kecepatan GEMM.

\section{Skenario Optimasi 3: Distributed-Memory (MPI)}

MPI diperlakukan sebagai eksperimen skalabilitas dan pembuktian jalur \keyword{out-of-core}, bukan sebagai \keyword{engine} yang ditukar pada demonstrasi antarmuka. Pada satu mesin dengan 4 \keyword{core}, empat proses MPI bersaing memperebutkan \keyword{core} dan \keyword{bandwidth} memori yang sama, sehingga \keyword{speedup} MPI tidak diharapkan melampaui OpenMP. Yang diukur adalah penurunan jejak memori per proses dan biaya komunikasi. Langkahnya:
\begin{enumerate}[leftmargin=2em]
  \item \keyword{Sharding}: indeks basis dibagi menjadi blok kontigu, dan tiap proses hanya memuat $\lceil N/P \rceil$ baris.
  \item \keyword{Broadcast}: proses utama menyiarkan matriks kueri yang berukuran kecil ke seluruh \keyword{worker} dengan \texttt{MPI\_Bcast}.
  \item \keyword{Gather dan merge}: tiap proses menghitung \keyword{top-k} lokal, lalu \texttt{MPI\_Gather} mengumpulkan kandidat ke proses utama untuk penggabungan eksak.

\end{enumerate}

% =====================================================================
\chapter{HASIL DAN EVALUASI}
% =====================================================================

\section{Parameter dan Kriteria Kelulusan}

Seluruh nilai pada bab ini diukur pada $N=100.000$, $D=384$, $Q=1.000$, $k=10$, dengan median dari 3 ulangan. Kriteria K7 diukur terpisah pada $N=500.000$. Kriteria kelulusan ditetapkan sebelum pengukuran agar tidak disesuaikan setelah hasil diketahui.

\begin{table}[H]
\centering
\caption{Kriteria kelulusan optimasi dan capaian.}
\label{tab:kpi}
{\small
\begin{tabularx}{\linewidth}{@{}c X c c l@{}}
\toprule
\textbf{Kode} & \textbf{Kriteria} & \textbf{Ambang} & \textbf{Hasil} & \textbf{Status}\\
\midrule
K1 & \keyword{Speedup} OpenMP 4 \keyword{thread} vs \keyword{baseline} & $\geq 2{,}5\times$ & $17{,}23\times$ & \pass\\
K2 & \keyword{Speedup} CUDA vs \keyword{baseline} (batch) & $\geq 20\times$ & $244{,}55\times$ & \pass\\
K3 & Latensi p95 CUDA per-kueri & $\leq 100$~ms & $2{,}59$~ms & \pass\\
K4 & Latensi p95 sistem \texttt{POST /search} & $\leq 100$~ms & $\approx 11{,}7$~ms & \pass\\
K5 & \keyword{Recall} top-$k$ tiap varian vs \keyword{baseline} & $=100\%$ & $100\%$ & \pass\\
K6 & \keyword{Parallel efficiency} OpenMP 4 \keyword{thread} & $\geq 62{,}5\%$ & $74{,}6\%$ & \pass\\
K7 & Jejak memori per proses MPI ($P=4$) & $\geq 3\times$ kecil & $3{,}70\times$ & \pass\\
\bottomrule
\end{tabularx}
}
\end{table}

\section{Throughput dan Speedup (Mode Batch)}

Tabel~\ref{tab:batch} menyajikan waktu eksekusi batch dan \keyword{speedup} relatif terhadap \keyword{baseline}.

\begin{table}[H]
\centering
\caption{Hasil mode batch, $N=100.000$, $D=384$, $Q=1.000$, $k=10$.}
\label{tab:batch}
\begin{tabular}{@{}l r r r r@{}}
\toprule
\textbf{Engine} & \textbf{Waktu (ms)} & \textbf{QPS} & \textbf{Speedup} & \keyword{Recall}\\
\midrule
\keyword{baseline} (1 thr) & 18.837,71 & 53,09 & $1{,}00\times$ & 100\%\\
\keyword{openmp} (1 thr)   & 3.262,45  & 306,52 & $5{,}77\times$ & 100\%\\
\keyword{openmp} (4 thr)   & 1.093,15  & 914,79 & $17{,}23\times$ & 100\%\\
\keyword{cuda}             & 77,03     & 12.982,15 & $244{,}55\times$ & 100\%\\
\bottomrule
\end{tabular}
\end{table}

Gambar~\ref{fig:speedup} menampilkan \keyword{speedup} pada skala logaritmik.

\begin{figure}[H]
\centering
\begin{tikzpicture}
\begin{axis}[
  ybar, bar width=22pt, width=0.82\linewidth, height=6.2cm,
  ymin=0.5, ymode=log,
  ylabel={\keyword{Speedup} vs baseline ($\times$)},
  symbolic x coords={openmp-1, openmp-4, cuda},
  xtick=data,
  xticklabels={OpenMP 1 thr, OpenMP 4 thr, CUDA (GEMM)},
  nodes near coords, point meta=explicit symbolic,
  enlarge x limits=0.28, grid=major,
]
\addplot+[ybar] coordinates {
  (openmp-1,5.77)   [5,77$\times$]
  (openmp-4,17.23)  [17,23$\times$]
  (cuda,244.55)     [244,55$\times$]
};
\end{axis}
\end{tikzpicture}
\caption{Speedup relatif terhadap baseline pada skala logaritmik.}
\label{fig:speedup}
\end{figure}

\section{Latensi Persentil Per-Kueri}

Untuk skenario interaktif, latensi diukur per kueri tunggal (Tabel~\ref{tab:latency}).

\begin{table}[H]
\centering
\caption{Latensi per-kueri (ms), $N=100.000$, $D=384$, $k=10$.}
\label{tab:latency}
\begin{tabular}{@{}l r r r@{}}
\toprule
\textbf{Engine} & \textbf{p50} & \textbf{p95} & \textbf{p99}\\
\midrule
\keyword{baseline} (CPU 1 thr) & 19,131 & 21,688 & 24,066\\
\keyword{cuda} (GEMM+refine)   & 1,475  & 2,586  & 3,071\\
\bottomrule
\end{tabular}
\end{table}

CUDA memenuhi K3 dengan p95 sebesar $2{,}586$~ms, jauh di bawah ambang 100~ms. Pada K4, latensi p95 endpoint \texttt{POST /search} dengan \keyword{engine} terbaik terukur sekitar $11{,}7$~ms pada $N=100.000$, juga memenuhi ambang.

\section{Skalabilitas dan Efisiensi Paralel}

\keyword{Speedup} paralel OpenMP diukur terhadap OpenMP 1 \keyword{thread} agar efek paralelisasi terpisah dari efek tiling dan SIMD. Hasilnya $3.262{,}45 / 1.093{,}15 = 2{,}98\times$ pada 4 \keyword{thread}, sehingga \keyword{parallel efficiency} $E = 2{,}98/4 = 74{,}6\%$, memenuhi K6 ($\geq 62{,}5\%$).

\section{Jejak Memori Terdistribusi (MPI)}

Karena tiap proses hanya memuat $\lceil N/P \rceil$ baris, jejak memori per proses menurun. Tabel~\ref{tab:mpi} menampilkan hasil pada $N=500.000$.

\begin{table}[H]
\centering
\caption{MPI: jejak memori per proses dan waktu komputasi ($N=500.000$, $Q=100$).}
\label{tab:mpi}
\begin{tabular}{@{}c r r@{}}
\toprule
\textbf{Proses ($P$)} & \textbf{Waktu komputasi (ms)} & \textbf{RSS/proses (KB)}\\
\midrule
1 & 9.349,07 & 770.532 ($\approx 752$~MB)\\
4 & 3.443,35 & 208.020 ($\approx 203$~MB)\\
\bottomrule
\end{tabular}
\end{table}

Reduksi jejak memori dari $P=1$ ke $P=4$ sebesar $770.532 / 208.020 = 3{,}70\times$, memenuhi K7 ($\geq 3\times$). Hasil MPI bersifat eksak, dengan himpunan \keyword{top-k} identik untuk $P=1,2,4$.

\section{Hasil Demonstrasi pada Korpus SIKA}

Pada aplikasi SIKA dengan korpus SOP terkurasi dan potongan referensi ($N=20.033$), \keyword{benchmark} 200 kueri per \keyword{engine} melalui \texttt{POST /bench} memberi hasil pada Tabel~\ref{tab:web}. CUDA menjadi \keyword{engine} dengan latensi terendah dan paling konsisten.

\begin{table}[H]
\centering
\caption{Latensi per-kueri pada korpus SIKA ($N=20.033$, $k=10$, 200 kueri).}
\label{tab:web}
\begin{tabular}{@{}l r r r r@{}}
\toprule
\textbf{Engine} & \textbf{p50 (ms)} & \textbf{p95 (ms)} & \textbf{p99 (ms)} & \textbf{QPS}\\
\midrule
\keyword{baseline} & 3,562 & 4,026 & 4,478 & 274,9\\
\keyword{openmp}   & 1,611 & 2,672 & 3,320 & 553,2\\
\keyword{cuda}     & 0,608 & 0,694 & 0,837 & 551,9\\
\bottomrule
\end{tabular}
\end{table}

Sebagai contoh \keyword{end-to-end}, pertanyaan ``Apakah pembelian lisensi perangkat lunak 40 juta perlu dua penawaran pembanding?'' dijawab dengan kutipan SOP-PGD-02 Pasal 1, latensi \keyword{retrieval} di bawah 1~ms, dan total respons termasuk LLM sekitar $2{,}5$~detik.

% =====================================================================
\chapter{ANALISIS BOTTLENECK}
% =====================================================================

\section{Memory-bound pada CPU}

Matriks basis $N=100.000$, $D=384$ berukuran sekitar $153{,}6$~MB, melampaui LLC 12~MB sekitar $12{,}8$ kali. \keyword{Baseline} membaca ulang seluruh matriks untuk setiap kueri. \keyword{Cache tiling} pada OpenMP meningkatkan pemakaian ulang baris basis, sehingga OpenMP 1 \keyword{thread} pun sudah $5{,}77\times$ lebih cepat dari \keyword{baseline} meskipun hanya memakai satu \keyword{thread}. Kenaikan dari 1 ke 4 \keyword{thread} menambah \keyword{speedup} menjadi $17{,}23\times$, sejalan dengan batas Hukum Amdahl.

\section{Compute-bound pada GPU}

GPU GTX 1650 dengan 4~GB VRAM unggul pada \keyword{throughput} masif. Dengan memetakan operasi inti ke GEMM cuBLAS dan memindahkan basis ke VRAM sekali, \keyword{speedup} batch mencapai $244{,}55\times$. Biaya pemindahan matriks sekitar 60~ms pada PCIe terbagi ke 1.000 kueri, sehingga tidak mendominasi pada mode batch.

\section{Komunikasi pada MPI}

Pada satu mesin 4 \keyword{core}, proses MPI bersaing memperebutkan \keyword{core} dan \keyword{bandwidth} memori. Karena itu MPI diposisikan sebagai jalur skalabilitas kapasitas (menurunkan jejak memori per proses), bukan sebagai percepatan latensi. Komunikasi \keyword{broadcast} untuk kueri berukuran kecil dan \keyword{gather} untuk \keyword{top-k} parsial menambah \keyword{overhead} yang terukur, namun tetap memberi penurunan waktu komputasi per proses.

% =====================================================================
\chapter{SISTEM TERDEMONSTRASI (SIKA)}
% =====================================================================

\section{Arsitektur Sistem}

SIKA dijalankan pada satu mesin dan dibagi menjadi empat lapisan:
\begin{enumerate}[leftmargin=2em]
  \item Lapisan Antarmuka (Web UI): Vite, React, dan TypeScript. Menampilkan jawaban, kutipan pasal, dan panel telemetri (p50/p95, QPS, \keyword{engine} aktif).
  \item Lapisan Layanan (API): FastAPI dengan endpoint \texttt{/ask}, \texttt{/search}, \texttt{/engine}, \texttt{/config}, \texttt{/models}, dan \texttt{/bench}.
  \item Lapisan Retrieval HPC: \keyword{engine} C++ (\keyword{baseline}, \keyword{openmp}, \keyword{cuda}, \keyword{mpi}) yang diekspos ke Python melalui \texttt{pybind11} dan dapat ditukar saat berjalan tanpa \keyword{restart}.
  \item Lapisan Data: matriks vektor \texttt{float32} $N \times D$ beserta metadata pasal.
\end{enumerate}

\section{Alur Permintaan \texttt{POST /ask}}

Pertanyaan diubah menjadi \keyword{embedding} (all-MiniLM-L6-v2, $D=384$), lalu \keyword{engine} aktif menjalankan $k$-NN, hasilnya dipetakan ke metadata pasal, potongan teratas disisipkan ke prompt LLM, dan jawaban dikembalikan bersama kutipan dan rincian latensi. Hanya tahap \keyword{retrieval} yang dioptimasi; tahap lain dijaga identik antar konfigurasi.

\section{Integrasi LLM}

LLM memakai protokol \keyword{OpenAI-compatible} yang dapat dikonfigurasi saat berjalan (Base URL, API Key, dan model melalui ``Fetch Models''). Bila penyedia gagal atau kena \keyword{rate-limit} (HTTP 429), sistem beralih ke jawaban \keyword{extractive} dari pasal hasil \keyword{retrieval}, sehingga demonstrasi tetap berjalan tanpa layanan eksternal.

\section{Skenario Demonstrasi}

Demonstrasi berdurasi sekitar delapan menit dengan alur berikut: D1 menjalankan sistem dan menampilkan status basis dokumen; D2 tanya normal dengan kutipan pasal; D3 beban 100 kueri serentak untuk menampilkan QPS dan p50/p95; D4 tukar \keyword{engine} secara langsung dari \keyword{baseline} ke \keyword{openmp} lalu \keyword{cuda}; D5 verifikasi \keyword{recall} 100\% antar \keyword{engine}; D6 jawab ulang pada \keyword{engine} tercepat di bawah satu detik; D7 jalur mundur CPU tanpa GPU.

% =====================================================================
\chapter{PENUTUP}
% =====================================================================

\section{Kesimpulan}

Akselerasi \keyword{exact} $k$-NN berlapis memenuhi seluruh kriteria kelulusan K1 sampai K7 dengan \keyword{recall} tetap 100\%. OpenMP memberi \keyword{speedup} $17{,}23\times$, CUDA memberi $244{,}55\times$ pada mode batch dengan p95 per-kueri $2{,}59$~ms, efisiensi paralel mencapai $74{,}6\%$, dan MPI menurunkan jejak memori per proses $3{,}70\times$. Optimasi tersebut terintegrasi ke dalam aplikasi SIKA yang dapat dioperasikan dan didemonstrasikan langsung, dengan \keyword{engine} yang dapat ditukar saat berjalan.

\section{Saran}

Pengembangan lanjutan dapat diarahkan pada implementasi hibrida MPI dan CUDA (multi-GPU per simpul), eksplorasi presisi rendah (\keyword{float16/bfloat16} dan \keyword{tensor core}), pembandingan terhadap pustaka industri seperti FAISS-GPU, serta penambahan korpus regulasi publik nyata untuk memperkuat validasi domain. Pada sisi aplikatif, purwarupa SIKA dapat dikembangkan menuju sistem multi-pengguna dengan autentikasi dan pembaruan indeks bertahap.

% =====================================================================
\appendix
\renewcommand{\thechapter}{\Alph{chapter}}
\renewcommand{\thesection}{\Alph{chapter}.\arabic{section}}
\setcounter{chapter}{0}
\setcounter{table}{0}
\setcounter{figure}{0}

\chapter{CARA REPRODUKSI}
% =====================================================================

\begin{lstlisting}[style=cppstyle,language=bash]
# 1) Build engine
cd engine && make bench bench_openmp bench_cuda mpi && cd ..
# 2) Binding Python (baseline + openmp + cuda)
bash scripts/build_binding.sh
# 3) Korpus demo (SOP terkurasi + potongan referensi)
python scripts/build_corpus.py --filler 20000
# 4) Evaluasi K1..K7 dan laporan
python scripts/evaluate_all.py --n 100000 --d 384 --q 1000 \
    --k 10 --reps 3 --mpi-n 500000
python scripts/make_report.py
# 5) Layanan dan web
bash scripts/run_service.sh            # http://localhost:8000
cd web && npm install && npm run dev   # http://localhost:5173
\end{lstlisting}

\noindent Untuk \keyword{engine} CUDA di WSL2, setel variabel lingkungan \texttt{LD\_LIBRARY\_PATH=/usr/lib/wsl/lib} agar \texttt{libcuda.so} milik WSL yang terpilih. Penyedia LLM yang berjalan di Windows diakses dari WSL melalui alamat IP host, bukan \texttt{127.0.0.1}; skrip \texttt{run\_service.sh} menangani hal ini secara otomatis.

\end{document}
